\input{preamble.tex}
\title{The smallest affine input in geometric twisting}
\author{Human Compression project}
\date{Private candidate, 11 October 2026}
\begin{document}
\maketitle
\begin{abstract}
Example 5.11 of Piedade--Tranchida, arXiv:2505.24662v1, says that a
particular affine quotient at parameter one is not residually connected,
and infers that its twist is not regular. The inference fails. We prove
from the literal presentations and coset subgroups that the twist is a
regular rank-three geometry for every positive parameter, including one.
No finite computation is needed for this conclusion.
\end{abstract}

\section{The source construction and its coordinates}
Let $W$ have the literal presentation on source page 45:
\[
 W=\langle a_1,a_2,a_3,a_4\mid a_i^2=1,\ [a_1,a_4]=[a_2,a_3]=1,
 (a_1a_2)^3=(a_1a_3)^3=(a_2a_4)^3=(a_3a_4)^3=1\rangle.
\]
Set $w=a_3a_4a_1a_2$, $N_s=\langle\!\langle w^{2s}\rangle\!\rangle$
and $A_s=W/N_s$, for $s\geq1$. Let $p$ and $q$ permute the four
named generators by $(12)(34)$ and $(23)$, respectively. They preserve
the defining Coxeter relations. Moreover,
\begin{equation}\label{eq:invariance}
 p(w)=(a_4a_3)w^{-1}(a_3a_4),\qquad
 q(w)=(a_3a_2)w(a_2a_3).
\end{equation}
The first identity is cancellation; the second also uses
$a_2a_3=a_3a_2$. Hence they preserve $N_s$ and induce an action
of the abstract diagram group
$B=\langle p,q\mid p^2=q^2=(pq)^4=1\rangle\cong D_8$
on $A_s$. The inverse permutations ensure equality of the normal
closures, rather than just one inclusion. Put $G_s=A_s\rtimes B$.
We use left-action coordinates
$(a,b)(c,d)=(a\,b(c),bd)$, so $bab^{-1}=b(a)$.
For the involutions $p,q$ this agrees with the source's conjugation
notation without a choice of inverse action.

In the source take $B_p=\langle q\rangle$,
$B_q=\langle p\rangle$, and choose $F_0=1$. The orbits of 1
under $1,B_p,B_q,B$ are, respectively,
$\{1\},\{1\},\{1,2\},\{1,2,3,4\}$. Thus its orbit-intersection
condition holds. The following are the literal three maximal parabolics of the twisting
system, as proved in Lemma~\ref{lem:literal} rather than inferred
from a quotient intersection property:
\begin{equation}\label{eq:maximal-literal}
 P_L=B,\qquad
 P_p=\langle a_1\rangle\rtimes\langle q\rangle,\qquad
 P_q=\langle a_1,a_2\rangle\rtimes\langle p\rangle.
\end{equation}
In particular these subgroups can be studied even when the input
coset system fails residual connectedness. We do not invoke a transfer
theorem whose input hypotheses are then absent.

\section{Identifying the literal quotient parabolics}
We retain the standard faithful affine Coxeter realization of the
literal presentation. This external input is needed here to identify
the actual quotient intersections in the source construction, as well
as for the order calculation and input-failure discussion below.
Let $Q=\{v\in\ZZ^4:\sum v_i=0\}$ and let $(\pi,v)$ act by
$z\mapsto\pi z+v$. In $W=Q\rtimes S_4$ take
\[
 a_1=((12),0),\ a_2=((23),0),\ a_3=((14),e_1-e_4),\ a_4=((34),0).
\]
Direct multiplication gives $w^2=(1,(1,1,-1,-1))$.
Consequently $N_s$ is the lattice $sM$, where
\[
 M=2Q+\ZZ(1,1,-1,-1)
  =\{v\in Q:v_i\equiv v_j\pmod2\text{ for every }i,j\}.
\]
Indeed conjugation permutes the coordinates of the translation;
differences of orbit vectors generate $2Q$, and every orbit vector
is congruent to $(1,1,-1,-1)$ modulo $2Q$. Thus $[Q:M]=4$,
$|A_s|=96s^3$ and $|G_s|=768s^3$. The affine model is an explicit retained premise of the literal-source
identification. The subsequent abstract marking argument uses only
the source presentation and an elementary $S_4$ quotient.

\begin{lemma}\label{lem:literal}
For every $s\geq1$, the actual quotient maximal parabolics satisfy
\[
 (A_s)_3\cap(A_s)_4=\langle a_1,a_2\rangle,\quad
 (A_s)_2\cap(A_s)_3\cap(A_s)_4=\langle a_1\rangle,\quad
 \bigcap_{i=1}^4(A_s)_i=1.
\]
Consequently the literal twisting subgroups are exactly
(\ref{eq:maximal-literal}).
\end{lemma}
\begin{proof}
The four maximal parabolics of $W$ fix the points
\[
 c_1=(3,-1,-1,-1)/4,\quad c_2=(1,1,-1,-1)/2,
 \quad c_3=0,\quad c_4=(1,1,1,-3)/4.
\]
Each has linear part $S_4$: its three generating transpositions form
an edge tree on the four coordinates. An element fixing $c_i$ has
translation $c_i-\pi c_i$. The quotient embeds each such parabolic,
since a pure translation fixing a point is the identity.

The third parabolic is the pure linear $S_4$. Its element $\pi$
is in the image of the fourth precisely when
$c_4-\pi c_4\in sM$. A nonzero such difference is a vector with
exactly two nonzero coordinates, one $1$ and one $-1$. It cannot
belong to $sM$: for $s\geq2$ divisibility fails, and for $s=1$
its four coordinate parities are not equal. Thus $\pi$ must fix
coordinate 4. These elements form exactly
$\langle(12),(23)\rangle=\langle a_1,a_2\rangle$.

Within this subgroup, membership in the second parabolic means
$c_2-\pi c_2\in sM$. Such a difference has all coordinates in
$\{-1,0,1\}$. It is zero exactly when $\pi$ preserves both
blocks $\{1,2\},\{3,4\}$. The only other way to have all
equal coordinate parities is to interchange the blocks, which is
impossible when coordinate 4 is fixed. Thus $\pi$ preserves the
blocks and fixes both 3 and 4, leaving only $\langle(12)\rangle$.
Finally membership in the first parabolic forces coordinate 1 to
be fixed, by the same two-odd-coordinates argument for
$c_1-\pi c_1$. This leaves the identity, proving all three formulas.
\end{proof}

\section{A uniform direct proof of regularity}
\begin{theorem}\label{thm:regular}
For every $s\geq1$, the literal coset system
$(G_s,(P_L,P_p,P_q))$ is a thin, residually connected,
flag-transitive rank-three geometry. Its distinguished marking is
$(x,p,q)$ with $x=a_1$, and its adjacent product orders are six
and four. In particular it is a regular hypertope at $s=1$.
\end{theorem}
\begin{proof}
The action of $B$ on the four names is transitive. Therefore $x,p,q$
generate $G_s$. The pair $a_1,a_2$ generates an exact $S_3$ in
every $A_s$. Its order is at most six by its two involutions and
braid of length three. For the lower bound use the quotient map
\[
 a_1\mapsto(12),\quad a_2\mapsto(23),\quad
 a_3\mapsto(14),\quad a_4\mapsto(34)
 \quad\text{in }S_4.
\]
All defining relations hold and the image of $w$ is $(12)(34)$,
so the map kills $w^{2s}$ for every $s$. The first two images
generate $S_3$, proving the lower bound and the claimed injection.
The same $S_4$ image distinguishes all four named generators in $A_s$.
The diagram group $D_8$ permutes these names faithfully, so its action
on $A_s$ is faithful as well. Thus our abstract semidirect factor is
exactly the source's actual group of diagram automorphisms.

Conjugation by $p$ interchanges $a_1,a_2$, whence
\[
 H_2=\langle x,p\rangle=\langle a_1,a_2\rangle\rtimes\langle p\rangle
 \quad\text{has order }12.
\]
Its involutory generators have product of order six:
$(xp)^2=a_1a_2$ has order three and the $B$ projection of $xp$
has order two. Also $H_0=\langle p,q\rangle=B$ has order eight,
and $H_1=\langle x,q\rangle\cong C_2^2$, since $q$ fixes $a_1$
and $x$ is nontrivial in the base. Semidirect coordinates give
\begin{equation}\label{eq:intersections}
 H_0\cap H_2=\langle p\rangle,\qquad
 H_0\cap H_1=\langle q\rangle,\qquad
 H_1\cap H_2=\langle x\rangle.
\end{equation}
For the last equality the projections lie in
$\langle q\rangle\cap\langle p\rangle=1$; the base part of
$H_1$ is exactly $\langle x\rangle$. The other two equalities
are immediate from the two-factor coordinates. All named generators
are distinct, so these identities include the complete intersection
property for the three-generator marking, with Borel subgroup one.
By (\ref{eq:maximal-literal}) these are precisely the source's maximal
parabolics, reordered as $(H_0,H_1,H_2)$.

We verify flag transitivity directly. Take three pairwise incident
cosets, one of each type. Translate a common representative of the
$H_0$ and $H_2$ cosets to the identity. The remaining coset $gH_1$
meets $H_0$ at some $u$ and $H_2$ at some $v$.
Then $v^{-1}u\in H_1$, so
$v^{-1}u=x^\epsilon q^\delta$. The element
$vx^\epsilon=uq^{-\delta}$ lies in all three cosets. Thus every
flag has a common representative, is a translate of a standard flag,
and extends to a chamber. Each rank-two residue has acting group one
of $H_0,H_1,H_2$, and its two child groups in
(\ref{eq:intersections}) generate it. Its incidence graph is connected:
a product of parabolic elements gives a path from the standard flag.
The whole incidence graph is connected because the maximal parabolics
generate $G_s$. This proves residual connectedness. Finally every
minimal parabolic has order two and the Borel is one, so all panels
have size two. The regular chamber action is faithful; an automorphism
fixing a chamber fixes each adjacent chamber in a thin geometry and
then all chambers by connectedness. Thus it is also the full
type-preserving automorphism group.
\end{proof}

\section{The input failure and the finite model}
The four maximal proper parabolics of $W$ fix, respectively,
\[
 c_1=(3,-1,-1,-1)/4,\quad c_2=(1,1,-1,-1)/2,
 \quad c_3=0,\quad c_4=(1,1,1,-3)/4.
\]
Their linear parts are $S_4$, and an element fixing $c_i$ has
translation $c_i-\pi c_i$. At $s=1$, the intersection of the
$c_2$ and $0$ stabilizers modulo $M$ is the order-eight group of
linear permutations preserving or interchanging the blocks
$\{1,2\},\{3,4\}$. This follows because $c_2-\pi c_2\in M$
exactly in those two cases. Intersecting additionally with the
$c_1$ stabilizer forces coordinate 1 to be fixed, leaving
$\langle(34)\rangle$; intersecting with the $c_4$ stabilizer
forces coordinate 4 to be fixed, leaving $\langle(12)\rangle$.
Nonzero differences $c_1-\pi c_1$ and $c_4-\pi c_4$ have exactly
two odd coordinates and hence cannot lie in $M$.
These two child groups generate only a group of order four, so the
rank-two residue of the flag of types 2 and 3 is disconnected.
The direct proof above shows that this input defect disappears in
the particular twisting construction.

The supplementary integer program constructs all four quotient maximal
parabolics and their actual intersections, and reconstructs the order-96 quotient,
the eight literal diagram automorphisms and their composition, and
the order-768 semidirect product. It compares all three literal
parabolics, all 64 generated-subgroup intersections and 343 proper
product identities. The full helper is retained alongside it; none
of its other finite parameter checks is a premise of the theorem.
There is no Lean claim. The exact source assertion being corrected
is the page-46 conclusion that the $s=1$ twist is not a regular
hypertope. This is a version-bound theorem/example correction, not a
numbered-conjecture solution, a full-source admission or a claim of
historical priority.
\end{document}
